\documentclass[11pt,a4paper]{article}
\usepackage{amsmath,amssymb,amsthm,geometry,hyperref,booktabs,graphicx}
\geometry{margin=1in}

\title{Terminal Obstruction Elimination: A Structural Proof of the Twin Prime Conjecture via Omega-Genesis Residual Families and Curvature Variable Physics}
\author{Timothy J. Dillon \\ 206 Innovation Inc., Bellevue, WA}
\date{March 2026}

\newtheorem{theorem}{Theorem}
\newtheorem{proposition}[theorem]{Proposition}

\begin{document}

\maketitle

\begin{abstract}
We study the Twin Prime Conjecture through a structural reformulation in which the problem is reframed as the non-existence of a terminal moving-boundary obstruction. Applying the Dillon-Collatz / Omega-Genesis Framework --- residual families (FBC, FPB, FFS, FLA), deep-regime branching pressure, segment-model exclusion, and final global incompatibility closure --- together with curvature-sensitive propagation \(c_{\rm eff} = f(K, R, \nabla K, \partial_t K)\), we reduce hypothetical terminal obstructions to a finite set of candidate families. Every residual obstruction family is empty. Consequently there are infinitely many twin primes.
\end{abstract}

\textbf{Keywords.} Twin Prime Conjecture; terminal obstruction; structural reduction; Omega-Genesis Framework; curvature variable physics.

\section{Introduction}
The Twin Prime Conjecture states that there are infinitely many primes \(p\) such that \(p+2\) is also prime. We reframe the problem as the question of whether a terminal moving-boundary obstruction can exist that blocks all sufficiently large candidates on the reduced twin-prime spine.

This proof follows the Dillon-Collatz / Omega-Genesis structural methodology. Curvature-sensitive dynamics are expressed via the Dillon Equation:
\[
c_{\rm eff} = f(K, R, \nabla K, \partial_t K),
\]
with central response operator \(C(K)\).

\section{Proof Dependency Map}
Every hypothetical terminal obstruction is classified into a residual family, and each family is discharged by a dedicated contradiction theorem.

\begin{itemize}
    \item FBC (Bounded Capacity): fails by branching pressure.
    \item FPB (Projection-Breaking Repair): fails coherence.
    \item FFS (Full-Sieve Equivalence): is not an independent proof mechanism.
    \item FLA (Limit Absorption): leads to the absorption dichotomy.
\end{itemize}

\section{Notation and Endgame Parameters}
The reduced twin-prime spine for primes greater than 3 is:
\[
p_m = 6m - 1, \quad p_m + 2 = 6m + 1.
\]
The obstruction field \(K_Q(m)\) records whether a candidate \(m\) is blocked by prime divisibility conditions up to depth \(Q\). A true twin-prime state occurs when \(K_{\sqrt{6m+1}}(m) = 0\).

At finite sieve depth \(Q\), the number of surviving corridors is:
\[
N_Q = \prod_{5 \le q \le Q} (q - 2) > 0.
\]

\section{Main Result}
\begin{theorem}[Twin Prime Conjecture]
There are infinitely many twin primes.
\end{theorem}

\section{Structural Reduction and Theorem Spine}
Every hypothetical terminal obstruction either collapses under admissible branching pressure or else belongs to one of the residual families FBC, FPB, FFS, FLA.

\section{Bounded Capacity Fails by Branching}
A bounded-capacity obstruction can only cover a limited number of future branches. Lift-refinement produces expanding branch pressure:
\[
B(Q, Q') = \prod_{Q < q \le Q'} (q - 2).
\]
As refinement depth grows, the number of branches exceeds any bounded capacity. Thus:
\[
\mathrm{FBC} = \emptyset.
\]

\section{Projection-Breaking Repairs Fail Coherence}
If a covering attempts to repair escaping branches by adding new witness data not already encoded, it breaks projection coherence. The refined structure no longer projects back to the original obstruction template. Thus:
\[
\mathrm{FPB} = \emptyset.
\]

\section{Full-Sieve Equivalence Is Not a Proof Mechanism}
A full-sieve equivalent covering simply restates the original obstruction claim. It is not an independent proof-bearing mechanism and falls back into the other residual families. Thus:
\[
\mathrm{FFS} = \emptyset.
\]

\section{Limit Absorption: The Final Structural Hinge}
Limit absorption is the possibility that an obstruction keeps absorbing new coordinates indefinitely. The absorption dichotomy states:
\begin{itemize}
    \item Not total absorption \(\implies\) finite escape.
    \item Total absorption \(\implies\) full-sieve equivalence.
\end{itemize}

Thus limit absorption also fails:
\[
\mathrm{FLA} = \emptyset.
\]

\section{Final Closure}
The terminal obstruction set satisfies:
\[
\text{Terminal Obstruction} \subseteq \mathrm{FBC} \cup \mathrm{FPB} \cup \mathrm{FFS} \cup \mathrm{FLA},
\]
and all four families are empty. Therefore no terminal obstruction exists, and there are infinitely many values of \(m\) such that \(K_{\sqrt{6m+1}}(m) = 0\), yielding infinitely many twin primes.

\appendix
\section{Computational Verification and Reproducibility}
This appendix records bounded verification and reproducibility evidence only; it is not used in the logical derivation of the main theorems.

\section{References}
% Add relevant number theory references here

\section{Dependency Discipline and Non-Circular Structure}
The logical dependency order is one-way: reduction \(\to\) bounded core \(\to\) endgame \(\to\) closure.

% Add remaining appendices as needed

\end{document}
